\begin{helper}[The Caccioppoli estimate (D.5) for the regularised signed power, at an arbitrary level and an arbitrary cut-off]
Let $n\ge2$, let $1<p<2$ and let $r>0$.  Let $C_1,C_2\in\R$ with
\[
C_1=\frac2p\,2^{1-\frac p2},
\qquad
C_2=\Bigl(\frac{2nrC_1p^2}{4(p-1)}\Bigr)^{2}.
\]
Let $\pi,\pi'\colon\R\times\R\to\R$ be two families of functions of one real
variable, written $\pi_\eta$ and $\pi'_\eta$, such that for every $\eta>0$:
\begin{enumerate}
\item[(H1)] $0<|s|+\eta$ for every $s\in\R$;
\item[(H2)] $(p-1)(|s|+\eta)^{p-2}\le\pi'_\eta(s)$ for every $s\in\R$;
\item[(H3)] $|\pi'_\eta(s)|\le\eta^{p-2}$ for every $s\in\R$;
\item[(H4)] $\pi'_\eta$ is continuous;
\item[(H5)] $\pi_\eta$ is $C^1$.
\end{enumerate}
Let $\varepsilon,\gamma_0,\gamma_1\in\R$, let
$A\colon\R^n\to\R^n$, let $u\colon\R^n\to\R$, let $g\colon\R^n\to\R^n$, let
$Dg\colon\{1,\dots,n\}\to(\R^n\to\R^n)$, let $x_0\in\R^n$, let $R\in\R$ and let
$k$ be an index.  Assume
\begin{enumerate}
\item[(H6)] $0<\varepsilon$, $0<\gamma_0$, $\gamma_1\le r\gamma_0$,
  $\gamma_0\le\gamma_1$ and $0<R$;
\item[(H7)] the field $(x,z,h)\mapsto A(h)$ satisfies \texttt{Q1} and
  \texttt{Q2} on $B_R(x_0)$ with the parameters
  $p,\varepsilon,\gamma_0,\gamma_1$ (\noderef{EllipticStructureQ12});
\item[(H8)] the cited regularity for $\varepsilon>0$ holds for $A$ on
  $B_R(x_0)$ with those parameters (\noderef{CitedRegularityEpsPos});
\item[(H9)] $u$ is a weak solution on $B_R(x_0)$ of the equation with field
  $(x,z,h)\mapsto A(h)$ and vanishing right-hand side, with weak gradient $g$
  (\noderef{IsWeakSolutionOn});
\item[(H10)] $g^i\in W^{1,2}_{\mathrm{loc}}(B_R(x_0))$ with weak gradient
  $Dg^i$, for every index $i$ (\noderef{MemW1pLocOn});
\item[(H11)] (almost-everywhere transport of a $W^{1,2}$ pair) for every
  $\Omega\subseteq\R^n$, all $f_1,f_2\colon\R^n\to\R$ and all
  $G_1,G_2\colon\R^n\to\R^n$: if $f_1\in W^{1,2}(\Omega)$ with weak gradient
  $G_1$ (\noderef{MemW1pOn}), $f_1=f_2$ almost everywhere on $\Omega$ and
  $G_1=G_2$ almost everywhere on $\Omega$, then $f_2\in W^{1,2}(\Omega)$ with
  weak gradient $G_2$;
\item[(H12)] for every $\eta>0$ and every $T$ with $0<T<R$, the function
  $\psi_\eta=\pi_\eta\circ g^k$ lies in $W^{1,2}(B_T(x_0))$ with weak gradient
  $\pi'_\eta(g^k)Dg^k$ (\noderef{MemW1pOn}).
\end{enumerate}
Write $w=\varepsilon+\|g\|^2$ and
$V_\eta=\tfrac p2(|g^k|+\eta)^{p/2-1}Dg^k$.  Then for every $\eta>0$ with
$\eta^2<\varepsilon$, every $\ell\in\R$ and every $\zeta\colon\R^n\to\R$ which
is $C^\infty$, compactly supported and satisfies
$\operatorname{supp}\zeta\subseteq B_R(x_0)$:
\begin{enumerate}
\item[(i)] the function $w^{\frac p2-1}\|V_\eta\|^2\zeta^2$ is integrable on
  $B_R(x_0)$;
\item[(ii)] with $A^+_\ell(\eta)=B_R(x_0)\cap\{\psi_\eta>\ell\}$,
  \[
  \int_{A^+_\ell(\eta)}w^{\frac p2-1}\|V_\eta\|^2\zeta^2\,dx
  \ \le\ C_2\int_{B_R(x_0)}\bigl((\psi_\eta-\ell)^+\bigr)^2|\nabla\zeta|^2\,dx ;
  \]
\item[(iii)] with $A^-_\ell(\eta)=B_R(x_0)\cap\{\psi_\eta<\ell\}$,
  \[
  \int_{A^-_\ell(\eta)}w^{\frac p2-1}\|V_\eta\|^2\zeta^2\,dx
  \ \le\ C_2\int_{B_R(x_0)}\bigl((\ell-\psi_\eta)^+\bigr)^2|\nabla\zeta|^2\,dx .
  \]
\end{enumerate}
Here $|\nabla\zeta|^2=\sum_i(D_i\zeta)^2$ and $(t)^+=\max\{t,0\}$.
\end{helper}
\begin{proof}
Throughout, $\eta>0$ with $\eta^2<\varepsilon$, $\ell\in\R$ and $\zeta$ are
fixed as in the statement, $B_R=B_R(x_0)$, and
$\psi_\eta=\pi_\eta\circ g^k$.  Note $\gamma_1>0$, since $0<\gamma_0\le\gamma_1$
by (H6).  For the consumer of this lemma $\pi_\eta(s)=(|s|+\eta)^{p-2}s$ and
$\pi'_\eta$ is its derivative, for which (H1)--(H5) are the thesis'
inequalities on p.~21 and (H12) is the chain rule for a $C^1$ function of a
$W^{1,2}$ function; but nothing below uses more than (H1)--(H5) and (H12).

\emph{Step 0: elementary consequences of the exponent range, and the weight.}
From $1<p<2$ and $n\ge2$ we get $p>0$, $p-1>0$, $p-2<0$ and $n>0$ as reals.
Since $\|g(x)\|^2\ge0$ we have $w(x)=\varepsilon+\|g(x)\|^2\ge\varepsilon>0$ at
every $x$, so $w^{s}$ is defined and positive for every real $s$; as
$p/2-1<0$, monotonicity of $t\mapsto t^{s}$ for $s<0$ gives
$w^{\frac p2-1}\le\varepsilon^{\frac p2-1}$ everywhere.  By (H8) there is a
continuous $g'$ on $B_R$ with $g=g'$ almost everywhere on $B_R$; hence
$w=\varepsilon+\|g'\|^2$ almost everywhere on $B_R$, the right-hand side is
continuous and bounded away from $0$ on $B_R$, and therefore for every real $s$
and every compact $K\subseteq B_R$ the function $w^{s}$ is almost everywhere
equal on $K$ to a continuous function and is bounded on $K$.  That is the only
use made of (H8): it supplies measurability and local boundedness of the
weight at the two exponents $p/2-1$ and $(p-2)/4$ that appear below.

\emph{Step 1: an intermediate radius.}  $\operatorname{supp}\zeta$ is a compact
subset of $B_R$, so there is $T$ with $0<T<R$ and
$\operatorname{supp}\zeta\subseteq B_T(x_0)$: if $\operatorname{supp}\zeta$ is
empty take $T=R/2$; otherwise the continuous function
$z\mapsto\operatorname{dist}(z,x_0)$ attains a maximum on it at some $y$, and
$T=(\operatorname{dist}(y,x_0)+R)/2$ works, since
$\operatorname{dist}(y,x_0)<R$.  The same construction gives such a $T$ for any
compact subset of $B_R$, and it is used in that generality below.

\emph{Step 2: the truncated test function.}  Put
\[
\theta^+_\eta=(\psi_\eta-\ell)^+,\qquad \varphi^+=\theta^+_\eta\,\zeta^2 .
\]
By (H12) at the radius $T$ of Step 1, $\psi_\eta\in W^{1,2}(B_T)$ with weak
gradient $\pi'_\eta(g^k)Dg^k$.  By \noderef{WeakGradientPositivePart} with
$q=2$, the level $c=\ell$, the open set $B_T$ of finite volume and that pair,
$\theta^+_\eta\in W^{1,2}(B_T)$ with weak gradient
$\chi_{\{\psi_\eta>\ell\}}\pi'_\eta(g^k)Dg^k$.  Since (H12) holds for
\emph{every} $T<R$ and the construction of $\theta^+_\eta$ does not involve
$\zeta$, and since every compact subset of $B_R$ sits inside some $B_T$ with
$T<R$ by Step 1, the same conclusion holds on every $B_T$, $T<R$; this is the
converse localisation, and it gives
$\theta^+_\eta\in W^{1,2}_{\mathrm{loc}}(B_R)$ with that weak gradient --- note
that it is cheaper here than in general because $\Omega$ is a ball, so the
intermediate radius is produced directly and no interposed compact set is
needed.  Now $\zeta^2$ is $C^\infty$ with compact support contained in $B_R$,
so clause (2) of \noderef{WeakGradientSmoothMul} with $q=2$ and $\Omega=B_R$
gives $\varphi^+\in W^{1,2}(B_R)$ with weak gradient
\[
g_{\varphi^+}=\zeta^2\chi_{\{\psi_\eta>\ell\}}\pi'_\eta(g^k)Dg^k
+\theta^+_\eta\,\nabla(\zeta^2),
\qquad \nabla(\zeta^2)=2\zeta\nabla\zeta ,
\]
and $\varphi^+$ has compact support with
$\operatorname{supp}\varphi^+\subseteq\operatorname{supp}\zeta\subseteq B_R$.
Where (H11) is needed is exactly here and in Step 1's instances: the pairs
produced by the composition and product rules agree almost everywhere with the
pairs written above, and (H11) transports membership along such an agreement.

\emph{Step 3: the differentiated equation against $\varphi^+$ (this is (D3)).}
Apply \noderef{DifferentiatedEquationWeak} with
$n,p,\varepsilon,\gamma_0,\gamma_1$ --- its hypotheses $2\le n$, $1<p$,
$0<\varepsilon$, $0<\gamma_0$ and $\gamma_0\le\gamma_1$ all being part of (H6)
--- the open set $B_R$, the field $A$, the solution $u$ with gradient $g$
(H9), the family $Dg$ (H10), the index $k$ and the pair
$(\varphi^+,g_{\varphi^+})$ of Step 2.  It gives
\[
\int_{B_R}\sum_{i,j}\bigl(DA(g)e_j\bigr)_i\,(Dg^k)_j\,(g_{\varphi^+})_i\,dx=0 .
\]
Separating the two summands of $g_{\varphi^+}$ --- legitimate because each is
integrable on $\operatorname{supp}\zeta$, by (H3), (H10) and Step 0, and both
vanish off $\operatorname{supp}\zeta$, so the split may be performed on
$\operatorname{supp}\zeta$ and then extended to $B_R$ --- and using that for
each fixed $i$ one has
$\sum_j\bigl(DA(g)e_j\bigr)_i(Dg^k)_j=\bigl(DA(g)Dg^k\bigr)_i$, because
$DA(g)$ is linear in its direction argument and $Dg^k=\sum_j(Dg^k)_je_j$, this
reads
\[
\int_{B_R}\chi_{\{\psi_\eta>\ell\}}\pi'_\eta(g^k)\,
\bigl\langle DA(g)Dg^k,Dg^k\bigr\rangle\zeta^2
=-2\int_{B_R}\theta^+_\eta\,\zeta\,
\bigl\langle DA(g)Dg^k,\nabla\zeta\bigr\rangle .
\]
On the left, \texttt{Q1} of (H7) with $\xi=Dg^k$ and $h=g(x)$ --- its guard
$\varepsilon+\|g\|^2=w\ne0$ holding since $\varepsilon>0$ --- gives
$\langle DA(g)Dg^k,Dg^k\rangle\ge\gamma_0w^{p/2-1}\|Dg^k\|^2$, and (H2) gives
$\pi'_\eta(g^k)\ge(p-1)(|g^k|+\eta)^{p-2}$, the two being combined because
both factors are nonnegative.  On the right, \texttt{Q2} gives
$|(DA(g)e_j)_i|\le\gamma_1w^{p/2-1}$, whence
\[
\bigl|\langle DA(g)Dg^k,\nabla\zeta\rangle\bigr|
\le\gamma_1w^{\frac p2-1}\sum_{i,j}|(Dg^k)_j|\,|(\nabla\zeta)_i|
\le n\gamma_1w^{\frac p2-1}\|Dg^k\|\,|\nabla\zeta| ,
\]
using $\sum_j|a_j|\le\sqrt n\,\|a\|$ (Cauchy--Schwarz, Preamble) twice.  Hence
\[
\gamma_0(p-1)\!\!\int_{A^+_\ell(\eta)}\!\!
w^{\frac p2-1}\|Dg^k\|^2(|g^k|+\eta)^{p-2}\zeta^2
\ \le\ 2n\gamma_1\!\!\int_{B_R}\!\!
w^{\frac p2-1}\|Dg^k\|\,\theta^+_\eta\,|\zeta|\,|\nabla\zeta| ,
\tag{D.3}
\]
the integrand on the right vanishing off $A^+_\ell(\eta)$ because
$\theta^+_\eta$ does.  This is the thesis' (3.13).

\emph{Step 4: the finiteness assertion (i).}  By (H1) and the definition of
$V_\eta$,
\[
\|V_\eta\|^2=\tfrac{p^2}4(|g^k|+\eta)^{p-2}\|Dg^k\|^2 ,
\]
which is the thesis' (3.14), and since $p-2<0$ and $|g^k|+\eta\ge\eta$ we get
$(|g^k|+\eta)^{p-2}\le\eta^{p-2}$, so
$\|V_\eta\|\le\tfrac p2\eta^{p/2-1}\|Dg^k\|$.  By (H10) and Step 1,
$\|Dg^k\|$ lies in $L^2$ of the compact set $\operatorname{supp}\zeta$;
$w^{p/2-1}\le\varepsilon^{p/2-1}$ by Step 0; and $\zeta$ is bounded and
vanishes off $\operatorname{supp}\zeta$.  Hence
$w^{\frac p2-1}\|V_\eta\|^2\zeta^2$ is integrable on $\operatorname{supp}\zeta$
and vanishes identically on $B_R\setminus\operatorname{supp}\zeta$, so it is
integrable on $B_R$, which is (i); measurability of the weight on
$\operatorname{supp}\zeta$ is Step 0 and measurability of $g^k$ and $Dg^k$
there is (H9)--(H10).  The same bound, applied on $A^+_\ell(\eta)$ in place of
$B_R$, shows that the left-hand side of (D.3) is finite.

\emph{Step 5: the Cauchy--Schwarz step (this is (D4)).}  Since
$\eta^2<\varepsilon$ we have $w\ge\varepsilon>\eta^2$, and $w\ge(g^k)^2$
because $\|g\|^2\ge(g^k)^2$; hence
$w^{1/2}\ge\max\{\eta,|g^k|\}\ge\tfrac12(|g^k|+\eta)$, and therefore, since
$1-p/2>0$,
\[
(|g^k|+\eta)^{1-\frac p2}\le\bigl(2w^{1/2}\bigr)^{1-\frac p2}
=2^{1-\frac p2}w^{\frac{2-p}4}.
\]
\emph{This is the only place where the hypothesis $\eta^2<\varepsilon$ is
used.}  From the identity of Step 4,
$\|Dg^k\|=\tfrac2p(|g^k|+\eta)^{1-\frac p2}\|V_\eta\|$, so
\[
w^{\frac p2-1}\|Dg^k\|\ \le\ \tfrac2p2^{1-\frac p2}\,
w^{\frac p2-1+\frac{2-p}4}\|V_\eta\|\ =\ C_1\,w^{\frac{p-2}4}\|V_\eta\| ,
\tag{D.4}
\]
because $\tfrac p2-1+\tfrac{2-p}4=\tfrac{p-2}4$.  Substituting (D.4) and the
identity $\|V_\eta\|^2=\tfrac{p^2}4(|g^k|+\eta)^{p-2}\|Dg^k\|^2$ into (D.3),
and applying the Cauchy--Schwarz inequality (Preamble) to the right-hand side
written as the product of $w^{\frac{p-2}4}\|V_\eta\|\,|\zeta|$ and
$\theta^+_\eta|\nabla\zeta|$ --- and noting
$\bigl(w^{\frac{p-2}4}\bigr)^2=w^{\frac p2-1}$, so that the first factor's
square integral is exactly the left-hand side ---
\[
\tfrac{4\gamma_0(p-1)}{p^2}\!\!\int_{A^+_\ell(\eta)}\!\!
w^{\frac p2-1}\|V_\eta\|^2\zeta^2
\ \le\ 2n\gamma_1C_1\Bigl(\int_{A^+_\ell(\eta)}\!\!
w^{\frac p2-1}\|V_\eta\|^2\zeta^2\Bigr)^{\frac12}
\Bigl(\int_{B_R}(\theta^+_\eta)^2|\nabla\zeta|^2\Bigr)^{\frac12}.
\]
Write $X$ for the integral on the left.  It is finite by Step 4 and
nonnegative.  If $X=0$ then (ii) holds outright, because its right-hand side is
nonnegative: $C_2\ge0$ and the integrand
$\bigl((\psi_\eta-\ell)^+\bigr)^2|\nabla\zeta|^2$ is nonnegative.  If $X>0$ we
may divide the displayed inequality by $X^{1/2}$ and square, obtaining
\[
X\ \le\ \Bigl(\frac{2n\gamma_1C_1p^2}{4\gamma_0(p-1)}\Bigr)^{2}
\int_{B_R}\bigl((\psi_\eta-\ell)^+\bigr)^2|\nabla\zeta|^2\,dx ,
\]
and $\gamma_1\le r\gamma_0$ with $\gamma_0>0$ turns the constant into
$\bigl(2nrC_1p^2/(4(p-1))\bigr)^2=C_2$.  That is (ii).

\emph{Step 6: the lower level sets (iii).}  By
\noderef{WeakGradientLinear} (clause (2), $q=2$, $a=-1$, $b=0$, both pairs
taken to be the pair of (H12)) the pair $(-\psi_\eta,-\pi'_\eta(g^k)Dg^k)$
lies in $W^{1,2}(B_T)$ for every $T<R$, and
\noderef{WeakGradientPositivePart} applied to it at the level $-\ell$ gives
$(\ell-\psi_\eta)^+\in W^{1,2}(B_T)$ with weak gradient
$-\chi_{\{\psi_\eta<\ell\}}\pi'_\eta(g^k)Dg^k$.  Steps 2--5 now run verbatim
with $\theta^-_\eta=(\ell-\psi_\eta)^+$ in place of $\theta^+_\eta$ and
$A^-_\ell(\eta)$ in place of $A^+_\ell(\eta)$: the sign of the test function's
weak gradient enters only inside the absolute values taken in Step 3, and the
lower bound of (H2) is applied on $A^-_\ell(\eta)$, where
$\chi_{\{\psi_\eta<\ell\}}=1$.  This gives (iii) and completes the proof.
\end{proof}
